\documentclass[11pt]{article}
\usepackage{mathblatt}
% Präambel des Setzers (werkzeuge/setzer-praeambel.tex, wortgleich); dazu \gteil (Muster T2U)
\makeatletter
\setlength{\mbpfbe}{0mm}% keine Punkte (pruefung.md § 5: nur beim Prüfstein)
\newlength{\gEin}\setlength{\gEin}{0pt}
\let\gSchrift\relax
\renewcommand{\pfkreuzzeile}[1]{\par\noindent\hangindent3.2em\hangafter1\hspace*{1.6em}$\square$\ #1\par}
\newcommand{\gkreuz}[1]{$\square$\ #1\hspace{2em}}
\newcommand{\gkreuzzeile}[2]{\par\noindent\hangindent2.2em\hangafter1\mbox{#1\,$\square$}\ #2\par}% Buchstabe, Option: eine je Zeile (neben Skizze)
% links, was man ansieht, rechts, was man tut; Spaltengrenze überall an derselben Stelle (gemeinsam.md „Satz“)
\newsavebox{\gbox}\newlength{\gLinks}
% v0.10 (Lehrer 13.10.): keine starre Halbierung – die linke Spalte ist so breit wie die Figur (+ 6 mm), die
% Antwort steht direkt daneben; ist die Figur breiter als 0,62 Zeile, steht die Antwort darunter
\newcommand{\gzwei}[2]{\par\smallskip\noindent\sbox{\gbox}{#1}%
  \setlength{\gLinks}{\dimexpr\wd\gbox+6mm\relax}%
  \ifdim\gLinks>0.62\linewidth
    \usebox{\gbox}\par\smallskip\noindent\begin{minipage}{\linewidth}#2\end{minipage}%
  \else
    \begin{minipage}[t]{\gLinks}\vspace{0pt}\usebox{\gbox}\end{minipage}%
    \begin{minipage}[t]{\dimexpr\linewidth-\gLinks\relax}\vspace{0pt}\raggedright #2\end{minipage}%
  \fi\par}
% v0.9 (12.10., pruefung.md § 5 „Marke mit B, Bedingungen rechts“): drittes Argument von pfaufg = Bedingungen
% der Prüfung (v0.10: \pfohnetr = TR durchgestrichen, „FOR“), klein grau im rechten Rand auf der Grundlinie der Nummer, Zeilen mit \\
\newcommand{\gbedrand}[2]{\ifblank{#2}{}{\rlap{\hspace*{\dimexpr#1+2mm\relax}%
  \parbox[t]{15mm}{\raggedright\scriptsize\leavevmode\color{mbgrau}\strut#2}}}}
% Teil eines Bündels (a, b …): Marke links im Rand, Bedingungen rechts im Rand, auf der Zeile des Teils
\newcommand{\gteilrand}[2]{\llap{\mbpfrandmarke{#1}\hspace*{\mbpfnr}}\gbedrand{\linewidth}{#2}}
% Teil eines Bündels: Buchstabe in eigener Spalte, Inhalt eingerückt (\gEin = Einzug, damit die Spaltengrenze von
% \gzwei an derselben Stelle bleibt)
\newenvironment{gbteil}[3]{\par\smallskip\noindent\gteilrand{#2}{#3}\makebox[\mbpfteil][l]{\textbf{#1}}%
  \begin{minipage}[t]{\dimexpr\linewidth-\mbpfteil\relax}\setlength{\gEin}{\mbpfteil}\raggedright\strut\ignorespaces}%
 {\end{minipage}\par}
% Grafik höchstens so breit wie die Spalte (Bündel nebeneinander, Koordinatensystem ohne Karo)
\newsavebox{\gpbox}
\newcommand{\gpassend}[1]{\sbox{\gpbox}{#1}\ifdim\wd\gpbox>\linewidth\typeout{GPASSEND: verkleinert \the\wd\gpbox\space auf \the\linewidth}\resizebox{\linewidth}{!}{\usebox{\gpbox}}\else\usebox{\gpbox}\fi}
% Skizzenfeld (v0.10, Lehrer 13.10.): Karo ab dem linken Rand über die Zeile, „Skizze“ klein grau oben links im Feld
% (wie \pfskizzenfeld der Vorlage); der rechte Rand bleibt den Bedingungen
\newcommand{\gskizze}[1]{\par\smallskip\noindent\begin{tikzpicture}
  \clip (0,0) rectangle ({\linewidth-1mm},{#1*10mm});
  \draw[black!16,line width=0.3pt,step=5mm] (0,0) grid ({\linewidth},{#1*10mm});
  \node[anchor=north west,inner sep=2pt,font=\scriptsize,text=mbgrau,fill=white] at (0,{#1*10mm}) {Skizze};
  \end{tikzpicture}\par}
\RenewDocumentEnvironment{pfaufg}{m m m}{%
  \begin{lrbox}{\mbaufgabebox}\begin{minipage}[t]{\linewidth}%
  \gSchrift\noindent\hspace*{\gEin}\mbpfrandmarke{#1}\gbedrand{\linewidth-\gEin}{#3}%
  \begin{minipage}[t]{\mbpfnr}\strut\textbf{#2}\end{minipage}%
  \begin{minipage}[t]{\dimexpr\linewidth-\gEin-\mbpfnr\relax}\raggedright\strut\ignorespaces}%
 {\end{minipage}%
  \end{minipage}\end{lrbox}%
  \setlength{\mbaufgabehoehe}{\dimexpr\ht\mbaufgabebox+\dp\mbaufgabebox\relax}%
  \ifdefempty{\gkopfnach}{\par\addvspace{5pt}\Needspace*{\mbaufgabehoehe}}{%
    % Stufenkopf, der mit dieser (ersten) Aufgabe zusammenbleibt: Platz für beide verlangen, dann den Kopf setzen
    % (12.10., Quadratische: ein Kopf stand allein unten, die Aufgabe mit Koordinatensystem auf der nächsten Seite)
    \par\Needspace*{\dimexpr\mbaufgabehoehe+40pt\relax}\gkopfnach\gdef\gkopfnach{}\par\addvspace{5pt}}%
  \noindent\usebox{\mbaufgabebox}\par\addvspace{5pt}}
% Stufenkopf vormerken: \gkopfvor{<Kopf>} – die nächste pfaufg setzt ihn, sobald beide zusammen Platz haben
\def\gkopfnach{}
\newcommand{\gkopfvor}[1]{\gdef\gkopfnach{#1}}
% Teilaufgabe mit Skizze links in gzwei: Buchstabe in eigener Spalte, darunter Text und Skizze
\newcommand{\gteil}[2]{\makebox[\mbpfteil][l]{\textbf{#1}}\parbox[t]{\dimexpr0.5\textwidth-\gEin-\mbpfnr-5mm-\mbpfteil\relax}{\raggedright #2}}
% Stufenkopf der Vorlage (mathblatt.sty 2026-10-12: Linie, Abstand, 13 pt, \pfjahre) ohne ihren Vorhalt von
% 14 Zeilen: der Setzer fordert vor dem Kopf 0,16 Seitenhöhe (Kopf + erste Aufgabe), mehr nicht
\setlength{\mbpfstufevorhalt}{0pt}
\makeatother
\newcommand{\karo}[1]{\par\smallskip\noindent\begin{tikzpicture}
  \clip (0,0) rectangle ({\linewidth-1mm},{#1*10mm});
  \draw[black!16,line width=0.3pt,step=5mm] (0,0) grid ({\linewidth},{#1*10mm});
  \end{tikzpicture}\par}
\newcommand{\kk}[1]{$\square$\,#1\hspace{0.9em}}
\newcommand{\linie}[1]{\,{\color{black!45}\rule[-1pt]{#1}{0.4pt}}\,}
\newcommand{\pkt}{\textperiodcentered{}}
% v0.14 (14.10.): Stufenkopf mit fünf Punkten (Option = Jahre für die Bestellkarte) oder „zuletzt 19“
\NewDocumentCommand{\gpunkte}{o m}{#2}
\newcommand{\gzuletzt}[1]{zuletzt #1}
% v0.7 (12.10.): Schülerfrage unter dem Stufenkopf; Ankreuzoption als Bild (Buchstabe und Kästchen darunter)
\newcommand{\gfrage}[1]{\par\nobreak\noindent\hspace*{\mbpfnr}{\small\itshape #1}\par\nobreak\addvspace{2pt}}
\newlength{\gbildbreite}\setlength{\gbildbreite}{0.32\linewidth}% Setzer setzt sie je Reihe (vier Bilder: 0,24)
\newcommand{\gbildoption}[3]{\begin{minipage}[t][#3][b]{\gbildbreite}\centering #2\par\smallskip\mbox{#1\,$\square$}\end{minipage}}
% Lineare (12.10.): Bild mit Buchstabe ohne Kästchen (Diagramm zuordnen); Aussage mit „wahr / falsch“ zum Ankreuzen
\newcommand{\gbildname}[3]{\begin{minipage}[t][#3][b]{\gbildbreite}\centering #2\par\smallskip\mbox{\textbf{#1}}\end{minipage}}
\newcommand{\gwahrfalsch}[2]{\par\noindent\hangindent2.2em\hangafter1\makebox[2.2em][l]{(#1)}#2\hfill\mbox{wahr $\square$\quad falsch $\square$}\par}
\makeatletter\setlength{\mbpfbe}{0mm}\makeatother
\begin{document}

\pfheftstil
\pflaufkopf[Abitur GK \pkt{} mehr]{Funktion aufstellen}{AA110}\pfseitevon
\pfheftkopf{Funktion aufstellen}{mehr}
\pfstufenzeile{Funktionsgleichung gesucht \pkt{} Zielfunktion aufstellen}
% 2018-be-gk-B1.1e-m1 (Schwere II, rang 6)
\begin{pfaufg}{GK~’18\\B1.1e}{1.}{}
Bei einer Sprungschanze wird die Anlaufbahn beschrieben durch $h(x) = 0{,}05x^2 + 30$ ($x \le 0$), der Aufsprunghang durch $g(x) = \tfrac{1}{1000} \cdot \left(\tfrac{1}{1600}x^4 - 8x^2 + 25\,600\right)$; 1 LE entspricht 1 m. Die Flugbahn ist Graph einer quadratischen Funktion $f$; in $S(0 \mid 30)$ geht die Anlaufbahn ohne Knick in die Flugbahn über. Bei $x = 40$ ist der Springer 4,4 m über dem Hang. Bestimme eine Funktionsgleichung von $f$.
\karo{8}
\fusshilfe{\mbox{1:~30; 0; $-$0,007x² + 30}}
\end{pfaufg}
% 2019-be-gk-B2.1c-m1 (Schwere II, rang 7)
\begin{pfaufg}{GK~’19\\B2.1c}{2.}{}
Die Höhe eines Baums in cm wird durch $g(t) = a \cdot t^3 + b \cdot t^2$ beschrieben; $t$ ist die Zeit in Jahren. Nach 4 Jahren ist der Baum 384 cm hoch und wächst dann 144 cm pro Jahr. Bestimme die Funktionsgleichung von $g$.
\karo{8}
\fusshilfe{\mbox{2:~3at² + 2bt; 144; $-$3t³ + 36t²}}
\end{pfaufg}
% 2021-be-gk-B2.1k-m1 (Schwere II, rang 8)
\begin{pfaufg}{GK~’21\\B2.1k}{3.}{}
Eine Parabel $q$ verläuft durch den Ursprung und hat den Hochpunkt $H(3 \mid 2{,}7)$. Bestimme eine Funktionsgleichung von $q$.
\karo{8}
\fusshilfe{\mbox{3:~0); 0; $-$0,3x² + 1,8x}}
\end{pfaufg}
% 2017-be-gk-B1.1f-m1 (Schwere II, rang 9)
\begin{pfaufg}{GK~’17\\B1.1f}{4.}{}
Ein Brückenteil einer Holzeisenbahn soll 20 cm lang sein, links 2 cm und rechts 10 cm hoch; in beiden oberen Eckpunkten soll die Profilkurve waagerecht verlaufen. Die linke untere Ecke liegt im Ursprung, 1 LE entspricht 1 cm. Das Profil wird durch $g$ mit $g(x) = ax^3 + bx^2 + c$ beschrieben. Bestimme die Funktionsgleichung von $g$.
\karo{8}
\fusshilfe{\mbox{4:~2; $-$30a; $-$0,002x³ + 0,06x² + 2}}
\end{pfaufg}
% 2018-be-gk-B1.1g-m1 (Schwere II, rang 10)
\begin{pfaufg}{GK~’18\\B1.1g}{5.}{}
Die Flugbahn eines Skispringers wird beschrieben durch $f(x) = -0{,}007x^2 + 30$, der Aufsprunghang durch $g(x) = \tfrac{1}{1000} \cdot \left(\tfrac{1}{1600}x^4 - 8x^2 + 25\,600\right)$; 1 LE entspricht 1 m. Weise nach, dass der senkrechte Abstand des Springers zum Hang höchstens 4,8 m beträgt. Auf die hinreichende Bedingung darfst du verzichten.
\karo{7}
\fusshilfe{\mbox{5:~0,001x² + 4,4 $-$ x$^{4}$/1 600 000; 28,3; Abstand höchstens 4,8 m}}
\end{pfaufg}
% 2021-be-gk-A1.3b-m1 (Schwere II, rang 11)
\begin{pfaufg}{GK~’21\\A1.3b}{6.}{}
Gegeben ist die Funktion $f$ mit $f(x) = x^3 - 6x^2 + 32$. Für $0 < x < 4$ legt der Punkt $P(x \mid f(x))$ zusammen mit dem Ursprung ein achsenparalleles Rechteck fest; sein Flächeninhalt ist für genau ein $x_{max}$ am größten. Weise nach, dass $x_{max} \ne 2$ ist.
\karo{6}
\fusshilfe{\mbox{6:~x$^{4}$ $-$ 6x³ + 32x; x\_max $\neq$ 2}}
\end{pfaufg}
% 2022-iqb-ga-A1.2a-m1 (Schwere I, rang 12)
\begin{pfaufg}{IQB~’22\\A1.2a}{7.}{}
Die Abbildung zeigt den Graphen einer Funktion $f$ mit $f(x) = a \cdot b^x$, $a, b > 0$. Bestimme die Werte von $a$ und $b$.
\par\smallskip\noindent\gpassend{\begin{abiksys}[xmin=-3,xmax=1.5,ymin=-0.5,ymax=3.5,le=0.8]\funktion[1.2]{1.5*2^\x}{}\punkt{0}{1.5}{}\punkt{1}{3}{}\end{abiksys}}\par
\karo{6}
\fusshilfe{\mbox{7:~1,5; 2}}
\end{pfaufg}
\par\vfill\noindent{\scriptsize\color{mbgrau}Originale: Abi GK 2018 \pkt{} B1.1e, Abi GK 2019 \pkt{} B2.1c, Abi GK 2021 \pkt{} B2.1k, Abi GK 2017 \pkt{} B1.1f, Abi GK 2018 \pkt{} B1.1g, Abi GK 2021 \pkt{} A1.3b, IQB GK 2022 \pkt{} A1.2a\par}
\end{document}